\enteteExam{NFP 107 -- Systèmes de gestion de bases de données - Durée 2h30}%
{Examen seconde session 21 avril 2021}%
{Effectué à distance}%


On souhaite mettre en place un système de prévention des maladies infectieuses
reposant sur le principe suivant\,: une application  
enregistre toutes les rencontres effectuées par le propriétaire
d'un smartphone avec d'autres personnes (équipées de la même application). 
Quand une infection par une maladie (comme la COVID19, la  varicelle, la rougeole, etc.)
est constatée chez une personne $P$, on doit pouvoir retrouver toutes  les
personnes $P_1, P_2, \cdots, P_n$ rencontrées pendant une période récente pour pouvoir
les informer.

Voici le schéma proposé. Les clés primaires sont en gras, les clés étrangères ne sont pas indiquées.

\begin{itemize}
\item Personne(\textbf{idPersonne},  prénom, nom, email, ville)
\item Rencontre (\textbf{idRencontre}, idPropriétaire, idPersonne, date)
\item Virus(\textbf{idVirus}, nom, duréeInfection)
\item Infection(\textbf{idPersonne}, \textbf{idVirus}, dateDébutInfection)
\end{itemize}

Dans la table \texttt{Rencontre}, \texttt{idPropriétaire} désigne le propriétaire
du smartphone qui enregistre la rencontre, et  \texttt{idPersonne} désigne
la personne rencontrée. Si, par exemple, Philippe (identifiant 11) est le,propriétaire
d'un smartphone équipé et rencontre Béatrice
(identifiant 56) à une date $d$, 
l'application insère dans la table 
\texttt{Rencontre} le nuplet 
$(id_x, 11, 56, d)$.  

\subsection*{Compréhension du schéma (6 pts)}

\begin{enumerate}
 \item Le schéma permet-il de représenter le fait qu'une personne a été infectée par plusieurs virus?

 \item Si je trouve le nuplet $(id_x, 11, 56, d)$ dans la base, cela signifie-t-il 
    que je vais trouver également un autre nuplet $(id_y, 56, 11, d)$\,? À votre avis
    dans quelles circonstances cette propriété peut être vérifiée.
  
 \item Donnez les commandes SQL de création des tables de ce schéma, 

  \item Quel est d'après vous le schéma entité-association 
        correspondant à ce schéma relationnel\,?
        
   \item Peut-on savoir avec ce schéma qu'une personne a été infectée plusieurs
      fois par le même virus\,? Expliquez, et proposez si nécessaire 
      une évolution du schéma pour que ce soit possible.

    \item Si je crée une table contenant l'union de tous les attributs des tables Virus et Infection,
        quelle serait la clé\,? Est-celle en troisième forme normale (donnez un argument précis)\,?  
\end{enumerate}

   
\begin{solution}
 \begin{enumerate}
    \item Oui, on peut avoir plusieurs fois la même
       personne mais avec des virus différents.
     \item Rien ne l'impose. Ce n'est vrai que si  
        les deux personnes qui se recontrent ont toutes
        les deux l'application.
     \item Standard. Voici pour la table \texttt{Rencontre}:
     
\begin{minted}{sql}
Create table Rencontre (idRencontre int not null, 
    idPropriétaire int not null, 
     idPersonne int not null, 
     date not null,
     primary key (idRencontre),
     foreign key (idPropriétaire) references Personne (idPersonne),
         foreign key (idPersonne) references Personne (idPersonne)
     )
\end{minted}
     
     \item 
     
     \item Non la clé interdit de stocker plusieurs
       fois la même paire (idPersonne, idVirus).
       Il vaut mieux réifier la table \texttt{Infection}
       en lui attribuant un identifiant propre.
     \item La clé est la paire (idPersonne, idVirus).
     Elle n'est pas en 3FN puisque le nom du virus
     ne dépend pas fonctionnellement de la clé 
     mais seulement d'une partie de la clé (la dépendance
     n'est donc pas minimale).
 \end{enumerate}
 
\end{solution}


\subsection*{SQL (7 pts)}

Sur le schéma obtenu précédemment, exprimez en SQL les requêtes suivantes\,:

\begin{enumerate}
  \item Donnez les nom et prénom des personnes infectées par le virus nommé  'varicelle'.

  \item Donnez les villes où on trouve  deux personnes
       infectées par le même virus (affichez la ville
       et les noms des deux personnes)

  \item Donnez la liste des personnes qui ont rencontré un propriétaire
      infecté par un virus APRES la date de début de l'infection
       (affichez le nom du propriétaire, le nom de la personne rencontrée
        et la date de rencontre)
   
   \item Donnez la liste des personnes qui n'ont pas été atteintes par le virus 'SARS'
  
   \item Donnez les propriétaires habitant Ajaccio et les personnes 
          qu'ils n'ont jamais rencontrées. Exemple: si Philippe habite Ajaccio
          et n'a jamais rencontré Béatrice, affichez ('Philippe', 'Béatrice").

   \item Affichez le nom de chaque virus avec le nombre de personnes
      infectées par ce virus
      
   \item Donnez la requête affichant le nom des propriétaires ayant rencontré au moins 1000 autres personnes
\end{enumerate}


\begin{solution}


\begin{enumerate}
  \item 
\begin{minted}{sql}
select prénom, nom
from  Personne as p, Infection as i, Virus as v
where p.idPersonne = i.idPersonne
and v.idVirus = i.idVirus
and v.nom='varicelle'
\end{minted}

 \item Il faut donc deux personnes (différentes mais
    dans la même ville),  deux infections pour le même virus.
\begin{minted}{sql}
select p1.ville,, p1.nom, p2.nom 
from  Personne as p1, Personne as p2, 
    Infection as i1, Infection as i2
where p1.idPersonne = i1.idPersonne
and p2.idPersonne = i2.idPersonne
and i1.idVirus = i2.idVirus
and p1.ville = p2.ville 
and p1.idPersonne < p2.idPersonne
\end{minted}

\item 
\begin{minted}{sql}
select p1.nom as nomPriopriétaire, p1.nom as nomPersonne,
     r.date
from  Rencontre as r, Personne as p1, Personne as p2, 
    Infection as i
where p1.idPersonne = r.idPropriétaire
and p2.idPersonne = r.idPersonne
and i.idPersonne =   p1.idPersonne
and r.date > i.dateDébutInfection
\end{minted}


\item 
\begin{minted}{sql}
select nom
from  Personne as p
where not exists (select * from Infection as i, Virus as v
     where i.idPersonne = p .idPersonne
     and i.idVirus = v.idVirus
     and v.nom='SARS')
\end{minted}


\item 
\begin{minted}{sql}
select p1.nom as nomPropriétaire, p2.nom as jamaisRencontrée
from  Personne as p1, Personne as p2
where not exists (select * from Rencontre as r
    where r.idPropriétaire = p1.idPersonne
    and r.idPersonne = p2.idPersonne)
and p1.ville='Ajaccio'
\end{minted}

\item 
\begin{minted}{sql}
select v.nom count(*) as nbInfections
from  Virus as v, Infection as i
where v.idVirus = i.idVirus
group by v.idVirus, v.nom
\end{minted}


\item 
\begin{minted}{sql}
select p.nom, Count(*) as nbRencontres 
from  Personne as p, Rencontre as r
where p.idPersonne = r.idPropriétaire
group p.idPersonne, p.nom
having count(*) >= 1000
\end{minted}
\end{enumerate}
\end{solution}


\subsection*{Algèbre (3 pts)}

On définit algébriquement les deux relations suivantes:

\begin{itemize}
  \item $P1 := \rho_{idPersonne \to id1, nom \to nom1, prenom \to prenom1, ville \to ville1} (Personne)$
  \item $P2 := \rho_{idPersonne \to id2, nom \to nom2, prenom \to prenom2, ville \to ville2} (Personne)$
\end{itemize}

Voici une expression algébrique:
$$\pi_{nom1, nom2} (\sigma_{ville1=ville2} (P1 \Join_{id1=idPropriétaire} Rencontre) \Join_{idPersonne=id2} P2)$$

\begin{enumerate}
  \item Expliquer en une phrase le but de cette requête
    \item Donnez la requête SQL correspondante
\end{enumerate}


\begin{solution}

\begin{enumerate}
  \item Les propriétaires et les personnes rencontrées qui habitent la même ville
  \begin{minted}{sql}
select prop.nom as nompropriétaire, pers.nom as nomPersonneRencontrée 
from  Personne as prop, Personne as pers, Rencontre as r
where prop.idPersonne = r.idPropriétaire
and pers.idPersonne = r.idPersonne
and prop.ville = pers.ville
\end{minted}
\end{enumerate}

\end{solution}

\subsection*{Procédures (2 pts)}

Quand un propriétaire $Prop$  a été infecté, on va rechercher toutes les personnes
en contact avec $Prop$ et leur envoyer un message. Pour ne pas envoyer une
notification deux fois, on ajoute un attribut ``notifié`` dans la table ``Rencontre``
qui vaut ``true``si une notification a été envoyée, `false``sinon. 
Voici la procédure appelée à chaque fois\,: elle prend comme paramètres le propriétaire
(personne infectée)
et la personne à notifier.

\begin{small}
\begin{minted}{sql}
create procedure Notification (v_prop RECORD, v_pers RECORD)  as 
begin
  -- Vérifions si la notification a été faite 
  select notifié into v_notifié from Rencontre as r
  where r.idPropriétaire = v_prop.idPersonne and r.idPersonne=v_pers.idPersonne;
  -- Pas de notification trouvée: on envoie un message et on mémorise l'envoi
  if (v_notifié == false) then
     update  Rencontre set notifié=true
     where  r.idPropriétaire = v_prop.idPersonne and r.idPersonne=v_pers.idPersonne;
     mail (v_pers.email, "Vous risquez d'avoir été infecté par " + v_prop.nom);
     commit;
   end ;
end;
\end{minted}
\end{small}

Expliquez pourquoi seul le mode ``serializable``garantit qu'on n'envoie pas une 
notification deux fois.



\begin{solution}
Si on lance deux fois la procédure au même moment pour les mêmes paramètres, 
chacune va conclure que la notification n'a pas eu lieu, et envoyer le message.

En mode sérialisable, cela ne peut pas arriver par définition: le résultat est identique
à celui obtenu par une exécution en série.
\end{solution}
\subsection*{Requêtes imbriquées(2 pts)}

Voici une requête avec imbrication.
\begin{small}
\begin{minted}{sql}
select distinct ville from Personne as p 
where  p.idPersonne in (select i.idPersonne from Infection as i
                        where i.idVirus in
                           (select v.idVirus from Virus as v
                            where v.nom = 'Peste'
                           )
                       )
\end{minted}               
\end{small}
Donnez une requête équivalente sans aucune imbrication.

\begin{solution}
\begin{minted}{sql}
select distinct ville 
from Personne as p, Infection as i, Virus as v
where  p.idPersonne = i.idPersonne 
and  i.idVirus  v.idVirus 
and  nom = 'Peste'
\end{minted}
\end{solution}
\end{document}

